\subsection{Solving the coupled equations}\label{e:solve}

The estimates proved so far concern solutions on the allowed side of
the floor.  To construct the compact map, we first solve the trace
equation for arbitrary bounded input $Z$; this step must not assume
$t\ge-\eps$.  Fix $n$, a sufficiently large finite truncation
$\Omega_R$, a homotopy parameter, and $0<b_{\rm reg}<1$.

\begin{lemma}[The scalar Dirichlet solve without a floor]\label{e:trace-solve}
For every invariant $Z\in C^{1,b_{\rm reg}}(\overline\Omega_R)$ with
outer value zero, the trace equation at the current homotopy stage has
a unique solution with its assigned Dirichlet data.  The solution
operator is continuous into $C^{2,b_{\rm reg}}$, also in the homotopy
parameter, and maps bounded input sets to bounded subsets of
$C^{3,b_{\rm reg}}$.  No lower bound for the implicitly defined $t$
is required.
\end{lemma}

\begin{proof}
We first work with smooth input $Z$ and obtain estimates depending
only on its $C^{1,b_{\rm reg}}$ norm.  Approximation will then give
the asserted solve for general inputs in that space.
Write the normalized trace equation as
\begin{equation}\label{e:normalized-trace}
 A_\chi:\nabla^2f=G_f,
 \qquad
 G_f=\frac{\sqrt D}{l}
           (F-\tr_{A_\chi}K).
\end{equation}
Here $F$ includes the modification appropriate to the current stage.
For an auxiliary parameter $s\in[0,1]$, interpolate between this
equation and $\Delta f=\tau_0f$:
\begin{equation}\label{e:scalar-interpolation}
 A^{(s)}:\nabla^2f=G^{(s)},\qquad
 A^{(s)}=(1-s)I+sA_\chi,\qquad
 G^{(s)}=(1-s)\tau_0f+sG_f.
\end{equation}
Keep the current boundary values throughout this interpolation.
At zero gradient $A_\chi=I$ and $\tr_{A_\chi}K=\tr K$.
The strict monotonicity of $F$ in $h$ therefore gives a uniform height
bound by the maximum principle, for every $s$ and every bounded input
set.

We next establish a global gradient bound for this scalar
interpolation.  If $Z$ is bounded and $\sigma\to\infty$, then
$t\to-\infty$.  In that range $p=n$ and $l=e^{nt}$, and the defining
relation between $t$ and $Z$ is exactly
\begin{equation}\label{e:nofloor-identity}
 e^{8Z}=e^{8t}+e^{(2n+8)t}\sigma^2.
\end{equation}
It follows, uniformly for bounded $Z$, that
\[
 d\asymp\sigma^{-8/(n+4)},\qquad
 \chi\asymp\sigma^{-8/(n+4)},\qquad
 \frac{\sqrt D}{l}\asymp\sigma.
\]
Because $n>4$, the exponent $8/(n+4)$ is less than one.  On the entire
bounded-height range we consequently have
\begin{equation}\label{e:nofloor-structure}
 \chi\ge\frac{c}{1+\sigma},\qquad
 |G_f|\le C(1+\sigma).
\end{equation}
The same inequalities hold for the longitudinal eigenvalue
$\chi^{(s)}=1-s+s\chi$ and the right side $G^{(s)}$, with uniform
constants.  The two transverse eigenvalues remain exactly one.

Let $r_0$ be distance from a boundary component in a fixed smooth
collar, and set
\begin{equation}\label{e:log-barrier}
 \psi(r_0)=\frac1k\log(1+B_0r_0),
 \qquad \psi''=-k(\psi')^2.
\end{equation}
Use $f_{\mathrm{face}}+\psi$ as an upper barrier and
$f_{\mathrm{face}}-\psi$ as a lower barrier.  At a possible contact their
gradient has length $q'=\psi'$, and the longitudinal second derivative
has magnitude $k(q')^2$.  By \eqref{e:nofloor-structure}, its signed
contribution dominates all forcing and collar-curvature terms:
\[
 \chi^{(s)}\psi''+O(q')
       \le-\frac{kc(q')^2}{1+q'}+Cq'.
\]
Choose $k$ large relative to the structural constants, then a collar
width $\delta'$ small enough that $1/(2k\delta')$ exceeds the required
slope threshold.  Finally choose $B_0\delta'\ge1$ so large that
$\psi(\delta')$ exceeds the height oscillation.  Throughout the collar
$\psi'\ge1/(2k\delta')$, so the differential inequalities have strict
sign; height monotonicity rules out a positive comparison maximum.
This proves
\begin{equation}\label{e:log-boundary}
 |f-f_{\mathrm{face}}|\le\psi(r_0)
                 \quad(0\le r_0\le\delta').
\end{equation}
The assigned values are constant on each face, and $\delta'$ can be
chosen below the separation of distinct boundary components.

For the interior estimate maximize
\[
 f(x)-f(y)-\psi(\dist(x,y))
\]
over pairs at distance at most $\delta'$.  The distance may be
computed in a fixed smooth extension of the metric; choose
$\delta'$ below its uniform injectivity radius.  Pairs with one point
on the boundary are controlled by \eqref{e:log-boundary}, since the
distance to that boundary component is no greater than the pair
distance.  Pairs at distance $\delta'$ are controlled by
$\psi(\delta')\ge\operatorname{osc}f$.  A positive maximum would
therefore occur at an interior pair with $r_0=\dist(x,y)>0$.
The two gradients there are parallel transports of the common vector
of length $q'=\psi'(r_0)$.  Simultaneous variations in parallel
transverse directions, using parallel trial fields in the second
variation of distance, give
\[
 \tr_\perp\nabla^2f(x)-\tr_\perp\nabla^2f(y)
                                      \le Cr_0q'.
\]
Separate longitudinal variations give
\[
 \nabla^2f(x)(e,e)\le\psi'',\qquad
 \nabla^2f(y)(e,e)\ge-\psi''.
\]
Subtracting \eqref{e:scalar-interpolation} at the two points yields
\begin{align*}
 -2C(1+q')
 &\le G^{(s)}(x)-G^{(s)}(y)\\
 &\le Cr_0q'+(\chi_x^{(s)}+\chi_y^{(s)})\psi''\\
 &\le Cr_0q'-\frac{2kc(q')^2}{1+q'}.
\end{align*}
Our choices of $k$ and the lower threshold for $q'$ make this
impossible.  Notice that the two longitudinal eigenvalues need not
coincide or have a known modulus of continuity.  Thus
$|f(x)-f(y)|\le\psi(\dist(x,y))$ for nearby pairs, proving a global
Lipschitz bound on each bounded input set.

The interpolation is now uniformly elliptic.  Lemma~\ref{g:gradient}
applies to its axial matrix and bounded forcing, first giving a
positive H\"older modulus for $\nabla f$.  Schauder estimates give
$C^{2,\theta}$ and then, by differentiating after this first gain,
$C^{2,b_{\rm reg}}$ and $C^{3,b_{\rm reg}}$ bounds.  More explicitly,
the first differentiated estimate gives $C^{3,\min(\theta,b_{\rm reg})}$;
hence $\nabla f$ is Lipschitz, allowing the original coefficients to
be estimated in $C^{b_{\rm reg}}$, followed by the stated full
exponent.  All these bounds are uniform in $s$.

For fixed input $Z$, the principal matrix is independent of the
height $f$, while
\[
 \partial_fG^{(s)}
 =(1-s)\tau_0+s\frac{\sqrt D}{l}\tau_0\partial_hF>0.
\]
The linearized Dirichlet operator thus has strictly negative
zeroth-order coefficient when written with everything on the left.
The maximum principle and linear elliptic Dirichlet theory make it
invertible \cite{GilbargTrudinger2001}.  Starting at
$\Delta f=\tau_0f$, openness follows from the implicit function
theorem, and the bounds above give closedness.  This proves existence
for $s=1$.  At a positive maximum of the difference of two solutions,
their gradients and principal matrices agree, while strict height
monotonicity contradicts the Hessian inequality; this proves
uniqueness.

For a general $C^{1,b_{\rm reg}}$ input, choose smooth invariant
inputs $Z_j$ with the same zero outer data, uniformly bounded in
$C^{1,b_{\rm reg}}$, and converging in $C^{1,b'}$ for every
$b'<b_{\rm reg}$.  Such an approximation follows by extension and
smoothing in boundary charts, correction of the outer trace, and
averaging under the circle action.  Convergence in the full
$C^{1,b_{\rm reg}}$ norm is neither asserted nor needed.  The smooth
scalar solutions have the uniform $C^{3,b_{\rm reg}}$ bounds just
proved.  Compact embedding gives a subsequence converging in
$C^{2,b_{\rm reg}}$ to a solution for $Z$; the third-derivative
H\"older seminorm bound passes to the limit by lower-exponent
compactness.  Uniqueness identifies this limit and removes dependence
on the approximation.  Finally, for any convergent sequence of
actual inputs in $C^{1,b_{\rm reg}}$, the same compactness and
uniqueness argument gives continuous $C^{2,b_{\rm reg}}$ dependence,
also at the matching endpoints of the four homotopy stages.
Symmetry follows from uniqueness.
\end{proof}

\begin{proposition}[Existence on the reduced exterior]\label{e:existence}
At each sufficiently large fixed $n$ there is a smooth invariant
physical solution of \eqref{d:trace}--\eqref{d:source} on the reduced
exterior, with its prescribed inner conditions and $t\ge-\eps$.
It is the local smooth limit of solutions on expanding truncations.
\end{proposition}

\begin{proof}
Use the Banach space $X_R$ of invariant
$C^{1,b_{\rm reg}}(\overline\Omega_R)$ functions with outer value
zero.  For each input $Z$, solve the trace equation by
Lemma~\ref{e:trace-solve}.  Compute $t,u,A_\chi,\Xi$ and all boundary
expressions from this input and its solved $f$.  Freeze them, replacing
only the differentiated $Z$ in the principal flux by the output
unknown $\widetilde Z$.  In the first three stages, for example, the
linear problem is
\[
 \divg\bigl(4uA_\chi\nabla\widetilde Z
               +uA_\chi\lambda K(w,\cdot)\bigr)=u\Xi,
\]
with the prescribed frozen inner flux and zero outer Dirichlet value.
The fourth stage uses its stated simultaneous source and boundary
scaling.

On bounded input sets the coefficients are uniformly elliptic.  The
nonempty outer Dirichlet face makes this mixed problem coercive;
there is no flux compatibility condition.  Its coefficients, drift,
and inner flux data are at least $C^{1,b_{\rm reg}}$, and its source
is at least $C^{0,b_{\rm reg}}$.  Linear boundary regularity therefore
bounds the output in $C^{2,b_{\rm reg}}$
\cite{AgmonDouglisNirenberg1959,GilbargTrudinger2001}.
Together with continuous dependence of the scalar solve, this defines
a continuous compact map $T_s:X_R\to X_R$, continuous in the full
homotopy parameter $s$.  The derivative count is worth specifying:
$f\in C^{3,b_{\rm reg}}$ and $Z\in C^{1,b_{\rm reg}}$ imply that
$t$, $uA_\chi$ and the drift are $C^{1,b_{\rm reg}}$, while the
quadratic expression $\mathcal T$ is $C^{0,b_{\rm reg}}$.
Uniqueness of both solves preserves invariance.  At a fixed point,
$Z\in C^{2,b_{\rm reg}}$ and $f\in C^{3,b_{\rm reg}}$ initially.
Then the frozen coefficients improve to $C^{2,b_{\rm reg}}$ and the
source to $C^{1,b_{\rm reg}}$, giving $Z\in C^{3,b_{\rm reg}}$;
the trace solve next gives $f\in C^{4,b_{\rm reg}}$.  Repeating this
triangular bootstrap makes each fixed point smooth.
Proposition~\ref{g:regularity} supplies the uniform classical bounds
for these smooth fixed points.

Let $t_s[Z]$ denote the value of $t$ after the scalar solve and define
the relatively open subset of $[0,1]\times X_R$
\[
 \mathcal O=\{(s,Z):t_s[Z]>-\eps
               \text{ on }\overline\Omega_R,\quad\|Z\|_{X_R}<M\}.
\]
Choose $M$ larger than the fixed-point bound supplied by
Proposition~\ref{e:ceiling} and Proposition~\ref{g:regularity}.
Lemma~\ref{d:floor} excludes fixed points on the floor boundary;
the classical estimates exclude them on the norm boundary.
Fixed points in the closure of $\mathcal O$ form a compact set,
because the homotopy maps bounded sets to relatively compact sets.
Leray--Schauder degree is therefore invariant on the moving sections
$\mathcal O_s$ \cite{LeraySchauder1934}.  One way to see the minor
moving-domain point is to cover this compact fixed-point set by
finitely many product neighborhoods lying in $\mathcal O$, subdivide
the parameter interval accordingly, and apply homotopy invariance and
excision on those fixed neighborhoods.  No control of $T_s$ on
unbounded input sets is needed.

At the terminal endpoint the trace comparison gives $f=0$, and the
linear second solve has identically zero output.  The zero input is
interior to $\mathcal O_1$, so its degree is one.  The physical
endpoint consequently has a fixed point with $t>-\eps$.
For this fixed $n$, let $R\to\infty$ through the admissible
truncations.  The estimates of Propositions~\ref{e:ceiling}
and~\ref{g:regularity} are uniform on every compact subset, including
inner-boundary charts.  A diagonal subsequence converges smoothly
there, giving the asserted physical solution and the limiting bound
$t\ge-\eps$.
\end{proof}

\subsection{End asymptotics and the energy flux}\label{e:flux-section}

\begin{lemma}[Asymptotics at fixed deformation parameter]\label{e:asymptotics}
At fixed $n$, the solution of Proposition~\ref{e:existence} has
exponentially decaying $f$, with derivatives through every required
finite order, and
\[
 Z,t=O_2(r^{-1}).
\]
In particular $\hat g=e^{4t}(g+l^2\dd f^2)$ is asymptotically flat of
order one.  Its energy is
\begin{equation}\label{e:energy-flux}
 E_{\hat g}=E_g-\frac{\mathfrak F_n}{8\pi},\qquad
 \mathfrak F_n=\lim_{R\to\infty}\int_{S_R}V_\nu
       =\int_{\Omega_{\rm red}}u\Xi+\int_BV_\nu.
\end{equation}
\end{lemma}

\begin{proof}
Beyond the support of $K,C,D_0,D_1$, the trace equation is
\[
 A_\chi:\nabla^2 f=\tau_0\frac{\sqrt D}{l}f.
\]
On the bounded variable range at fixed $n$ its principal matrix is
uniformly elliptic and its height coefficient has a positive lower
bound.  Slowly decaying exponentials $Ce^{-c(r-r_*)}$ are positive
supersolutions for the operator obtained by moving the height term to
the left, if $c>0$ is small.  Choose $C$ to dominate the data on
$r=r_*$.  The zero outer data allow the same barriers on every
truncation; applying them to both signs and then exhausting gives
exponential decay of $f$.  For the derivative assertion, freeze the
coefficients in this homogeneous linear equation along the solution.
Proposition~\ref{g:regularity} bounds their $C^m$ norms on uniform
unit patches, independently of $R$, for every fixed finite $m$.
Local homogeneous estimates therefore give
$\|f\|_{C^{j,\alpha}(B_1)}\le C(n,j)\|f\|_{C^0(B_2)}$;
on outer half-patches the same estimate uses the zero Dirichlet
condition.  Each finite $j$ requires only finitely many of the already
bounded coefficient derivatives.  Thus every required derivative of
$f$ decays exponentially; no decay of the derivatives of $Z$ is
needed at this step.

For clarity, put $E=t-Z=-\tfrac18\log D$.  Then $E$, $D-1$,
$v$, $A_\chi-I$, and their required derivatives are
$O_n(e^{-2cr})$, while $w,H^f,F=O_n(e^{-cr})$ with derivatives.
The tensor expression for $\mathcal T$ consequently gives
\[
 \mathcal T=4[p(Z)+1]|\dd Z|^2+O_n(e^{-2cr}).
\]
Every remaining term contains two factors from $H^f,F,w$, or a graph
error multiplying bounded derivatives of $Z$.  In particular,
$u-L(Z)$ and $p(t)-p(Z)$ have the same exponential bound, and
\[
 V=4L(Z)\nabla Z+E_V,\qquad
 \divg E_V=O_n(e^{-2cr}).
\]
Here the divergence uses only the radius-uniform unit-patch bound on
$\nabla^2Z$.  The height term $\tau_0a\sigma$ and the penalty term
$v\rho_1$ are exponentially small as well.  There are no persistent
$K$ terms, since $K$ vanishes on this end; derivatives of the profiles
have bounded fixed-$n$ coefficients and multiply the stated graph
errors.  Consequently the end divergence equation takes the form
\begin{equation}\label{e:end-Z-equation}
 \Delta Z+
 \bigl[p(Z)+2-\delta_0(p(Z)+1)\bigr]|\nabla Z|^2
                                     =O(r^{-4}).
\end{equation}
This also holds on the truncations, with uniform constants.  More
precisely, the right side is
\[
 \frac14\rho-\frac14C_nm_0(Z)\rho_0+O_n(e^{-2cr}),
\]
which is $O(r^{-4})$ without any prior decay estimate for $Z$.

Put $b(z)=p(z)+2-\delta_0(p(z)+1)$ and choose $\Phi(0)=0$,
$\Phi'(z)=\exp(\int_0^z b(q)\,\dd q)$.  On the bounded solution
range, $\Phi'$ and its reciprocal are bounded.  Then
$Y=\Phi(Z)$ satisfies $\Delta Y=O(r^{-4})$.  For
$0<\eta<1$, the positive function
$r^{-1}-r^{-1-\eta}$ has Laplacian
$-\eta(1+\eta)r^{-3-\eta}+O(r^{-4})$ in the AF metric.
Multiples of this function dominate both the source and the bounded
inner-end data, and are positive at the zero outer Dirichlet face.
Comparison on the truncations proves $Y=O(r^{-1})$, hence
$Z=O(r^{-1})$ after exhaustion.

The penalty therefore vanishes sufficiently far out.  In the
transformed equation the remaining nonexponential terms have smooth
weight coefficients of order $-4$ and bounded smooth factors
in $Z$.  On rescaled annuli, a $W^{2,p}$ estimate with $p>3$ first
bounds the rescaled gradient and its H\"older modulus.  Schauder
estimates then bound the rescaled second derivatives.  The
exponentially small graph errors are controlled in these estimates
by the unit-patch bounds of every required finite order.  This yields
$Z,t=O_2(r^{-1})$.

It follows that $\divg V=u\Xi$ is integrable and
$V=4\nabla t+O(r^{-3})$.  The divergence theorem on $\Omega_R$
gives
\[
 \int_{S_R}V_\nu=\int_{\Omega_R}u\Xi+\int_BV_\nu,
\]
because the outward normal of $\Omega_R$ on $B$ is $-\nu$.
This proves existence of the flux limit.  The graph term has
exponentially decaying ADM contribution.  For the conformal factor
$e^{4t}$ the ADM change is
$-(2\pi)^{-1}\lim\int_{S_R}\partial_\nu t$, which is precisely
$-\mathfrak F_n/(8\pi)$.  This proves \eqref{e:energy-flux}.
\end{proof}

\begin{proposition}[Vanishing negative flux]\label{e:flux}
With the prepared data and $\eps>0$ fixed,
\[
 \mathfrak F_n\ge-o_n(1),\qquad
 E_{\hat g}\le E_g+o_n(1)
                    \quad\text{as }n\to\infty.
\]
\end{proposition}

\begin{proof}
At the physical endpoint, \eqref{e:small-boundary-flux} bounds the
negative inner flux by $C(\tau_0/\ell)\int_B L$.
Apply \eqref{e:weighted-trace-eq} with a cutoff equal to one near $B$.
Since $\rho$ is bounded below on the collars,
\[
 1+\sqrt{\mathcal T}\le C(\rho+\delta_0\mathcal T)
\]
there.  As $\Pi_n\tau_0/\ell\to0$, this cost absorbs into, for
example, half of the positive bulk integral of
$u(\rho+\delta_0\mathcal T)$.  These are polynomial estimates; none
of the arbitrary fixed-$n$ regularity constants enters this
absorption.

On the penalty support, $-\eps\le t\le-\eps+1/n$, so
$l$ and $L$ are comparable to $\ell$.  Directly from the definitions,
\begin{equation}\label{e:penalty-sizes}
 u\le C(\ell+\ell^2\sigma),\qquad
 uv\le C\ell^2\sigma,\qquad
 ua\sigma=Ll\sigma^2\ge c\ell^2\sigma^2.
\end{equation}
Therefore
\[
 um_0\mathcal P\le
       \Pi_n\bigl[\ell\rho_0+
                        \ell^2\sigma(\rho_0+\rho_1)\bigr].
\]
Young's inequality against the remaining positive
$\delta_0\tau_0ua\sigma$ term yields
\begin{equation}\label{e:penalty-absorption}
 um_0\mathcal P\le
 \frac{\delta_0\tau_0}{2}ua\sigma+
 \Pi_n\left[\ell\rho_0+
       \frac{\ell^2}{\tau_0}(\rho_0^2+\rho_1^2)\right].
\end{equation}
The displayed remainder has vanishing total integral.
Indeed $\rho_0\asymp r^{-4}$ is integrable,
$\rho_0^2$ is integrable, and
$\rho_1^2\asymp r^{-4\gamma}$ is integrable in dimension three
because $\gamma>3/4$.  Finally
\[
 \ell=e^{-n\eps},\qquad
 \frac{\tau_0}{\ell}=\ell^{1/2},\qquad
 \frac{\ell^2}{\tau_0}=\ell^{1/2},
\]
so both small factors dominate every polynomial loss in $n$.
Insert these bounds in \eqref{e:energy-flux} to obtain the proposition.
\end{proof}

\begin{proof}[Completion of the proof of Theorem~\ref{d:deformation}]
At the physical endpoint $F=\tau_0f+C$, whence
$w(F)=\tau_0a\sigma+w(C)$.  Substituting \eqref{d:source} into
\eqref{d:curvature} shows that the remainder after
$8\pi(\mu+J_0(w))$ is
\[
 (1-\delta_0)\mathcal T+
 (1-\delta_0)\tau_0a\sigma+m_0\mathcal P
                     +w(C)-F\tr K-\rho.
\]
The first three terms are nonnegative, and the height and offset
bounds give
\[
 w(C)-F\tr K-\rho
 \ge-|\nabla C|-|F\tr K|-\rho\ge-\mathfrak m/2.
\]
This is the scalar-curvature conclusion of the theorem.
Combining the physical boundary condition with
\eqref{d:boundary-identity} gives
$H+4\partial_\nu t=-N_B$, so
\[
 H_{\hat g}=-e^{-2t}\sqrt d\,N_B<0.
\]
The metric inequality $\hat g\ge e^{-4\eps}g$ guarantees
completeness with the inner boundary included.
Lemma~\ref{e:asymptotics} supplies the required end regularity and
finite energy, and Proposition~\ref{e:flux} supplies its upper bound.
All choices respect axial invariance.  The limits are taken in the
order established above: fix the prepared data, selected domains,
cutoff shapes and $\eps$; choose large $n$ and form its profiles;
solve and exhaust $R$ at that fixed $n$; then let $n\to\infty$ in
the scalar energy estimate.
No convergence of the deformation metrics as $n\to\infty$ is needed.
\end{proof}
